\documentclass[a4paper]{article}

\usepackage[margin=1in]{geometry}
\usepackage{amsmath,amsfonts,mathtools,amsthm,amssymb}
\usepackage[l2tabu,orthodox]{nag}
\usepackage{microtype} % fixes spacing or whatever
\usepackage{enumitem}
\usepackage{tikz-cd}
\usepackage{xcolor}
\usepackage{physics}
\usepackage{mathrsfs}
\usepackage{cancel}
% \usepackage{libertine,libertinust1math}
\usepackage{eucal}
\usepackage[toc,page]{appendix}
\usepackage{pgfplots}
\pgfplotsset{compat=1.18}

\newtheorem{proposition}{Proposition}
\newtheorem{theorem}{Theorem}
\newtheorem{lemma}{Lemma}
\newtheorem{corollary}{Corollary}

\theoremstyle{definition}
\newtheorem{remark}{Remark}
\newtheorem{definition}{Definition}
\newtheorem{example}{Example}

% \usepackage[sc,osf]{mathpazo}   % With old-style figures and real smallcaps.
% \linespread{1.025}              % Palatino leads a little more leading
% % Euler for math and numbers
% \usepackage[euler-digits,small]{eulervm}
% \AtBeginDocument{\renewcommand{\hbar}{\hslash}}

\usepackage{etoolbox}
\usepackage{dynkin-diagrams}
\def\row#1/#2!{\dynkin{#1}{#2}\\}
\newcommand{\tble}[1]{
   \renewcommand*\do[1]{\row##1!}
   \[
      \begin{array}{ll}\docsvlist{#1}\end{array}
   \]
}

\allowdisplaybreaks

\renewcommand{\epsilon}{\varepsilon}
% \renewcommand{\phi}{\varphi}

\DeclareMathAlphabet{\mathpzc}{OT1}{pzc}{m}{it}

\newcommand{\goth}[1]{\mathfrak{#1}}
\newcommand{\red}[1]{#1_{\text{red}}}
\newcommand{\ad}[1]{#1_{\text{Ad}}}
\newcommand{\reg}[1]{#1_{\text{reg}}}
\newcommand{\sh}[1]{#1^{\text{sh}}}
\newcommand{\cpdv}[2]{\cfrac{\partial #1}{\partial #2}}

\newcommand{\fA}{\mathcal{A}}
\newcommand{\fB}{\mathcal{B}}
\newcommand{\fC}{\mathcal{C}}
\newcommand{\fD}{\mathcal{D}}
\newcommand{\fE}{\mathcal{E}}
\newcommand{\fF}{\mathcal{F}}
\newcommand{\fG}{\mathcal{G}}
\newcommand{\fH}{\mathcal{H}}
\newcommand{\fI}{\mathcal{I}}
\newcommand{\fJ}{\mathcal{J}}
\newcommand{\fK}{\mathcal{K}}
\newcommand{\fL}{\mathcal{L}}
\newcommand{\fM}{\mathcal{M}}
\newcommand{\fN}{\mathcal{N}}
\newcommand{\fO}{\mathcal{O}}
\newcommand{\fP}{\mathcal{P}}
\newcommand{\fQ}{\mathcal{Q}}
\newcommand{\fR}{\mathcal{R}}
\newcommand{\fS}{\mathcal{S}}
\newcommand{\fT}{\mathcal{T}}
\newcommand{\fX}{\mathcal{X}}
\newcommand{\fY}{\mathcal{Y}}
\newcommand{\fZ}{\mathcal{Z}}

\newcommand{\be}{\mathbf{e}}
\newcommand{\bx}{\mathbf{x}}
\newcommand{\bone}{\mathbf{1}}
\newcommand{\cx}{\bar{x}}
\newcommand{\cy}{\bar{y}}
\newcommand{\vx}{\vec{x}}
\newcommand{\vsigma}{\vec{\sigma}}
\newcommand{\tzeta}{\tilde{\zeta}}

\newcommand{\gm}{\goth{m}}
\newcommand{\gp}{\goth{p}}
\newcommand{\gF}{\goth{F}}
\newcommand{\gU}{\goth{U}}
\newcommand{\gV}{\goth{V}}

\newcommand{\A}{\mathbf{A}}
\newcommand{\C}{\mathbf{C}}
\newcommand{\F}{\mathbf{F}}
\newcommand{\G}{\mathbf{G}}
\newcommand{\bH}{\mathbf{H}}
\newcommand{\N}{\mathbf{N}}
\newcommand{\PP}{\mathbf{P}}
\newcommand{\Q}{\mathbf{Q}}
\newcommand{\R}{\mathbf{R}}
\newcommand{\Z}{\mathbf{Z}}

\newcommand{\dlog}{\dd\log}

\newcommand\srestr[2]{{\left.\kern-\nulldelimiterspace #1\vphantom{\small|} \right|_{#2}}}
\newcommand\restr[2]{{\left.\kern-\nulldelimiterspace #1 \vphantom{\big|} \right|_{#2}}}
\newcommand\gsec[2]{\Gamma({#1}, {#2})}
\newcommand\gsecx[1]{\Gamma(X, {#1})}

\DeclareMathOperator{\chr}{char}
\DeclareMathOperator{\rProj}{\mathbf{Proj}}
\DeclareMathOperator{\rSpec}{\mathbf{Spec}}
\DeclareMathOperator{\rHom}{\mathbf{Hom}}
\DeclareMathOperator{\rExt}{\mathbf{Ext}}
\DeclareMathOperator{\id}{id}
\DeclareMathOperator{\Frac}{Frac}
\DeclareMathOperator{\rk}{rank}
\DeclareMathOperator{\deter}{det}
\DeclareMathOperator{\pic}{Pic}
\DeclareMathOperator{\cacl}{CaCl}
\DeclareMathOperator{\trd}{tr.d.}
\DeclareMathOperator{\cl}{Cl}
\DeclareMathOperator{\depth}{depth}
\DeclareMathOperator{\codim}{codim}
\DeclareMathOperator{\Div}{Div}
\DeclareMathOperator{\coker}{coker}
\DeclareMathOperator{\len}{length}
\DeclareMathOperator{\height}{height}
\DeclareMathOperator{\supp}{Supp}
\DeclareMathOperator{\proj}{Proj}
\DeclareMathOperator{\im}{im}
\DeclareMathOperator{\Hom}{Hom}
\DeclareMathOperator{\Der}{Der}
\DeclareMathOperator{\spec}{Spec}
\DeclareMathOperator{\Aut}{Aut}
\DeclareMathOperator{\ch}{char}
\DeclareMathOperator{\tor}{Tor}
\DeclareMathOperator{\Ann}{Ann}
\DeclareMathOperator{\Syl}{Syl}
\DeclareMathOperator{\Sym}{Sym}
\DeclareMathOperator{\GL}{GL}
\DeclareMathOperator{\SL}{SL}
\DeclareMathOperator{\Stab}{Stab}
\DeclareMathOperator{\Perm}{Perm}
\DeclareMathOperator{\Orb}{Orb}
\DeclareMathOperator{\Gal}{Gal}
\DeclareMathOperator{\Supp}{Supp}
\DeclareMathOperator{\Ext}{Ext}
\DeclareMathOperator{\Tor}{Tor}
\DeclareMathOperator{\PGL}{PGL}
\DeclareMathOperator{\hd}{hd}
\DeclareMathOperator{\pd}{pd}
\DeclareMathOperator{\SO}{SO}
\DeclareMathOperator{\SU}{SU}
\DeclareMathOperator{\Sp}{Sp}
\DeclareMathOperator{\Oth}{O}
\DeclareMathOperator{\Uni}{U}
\DeclareMathOperator{\evf}{ev}
\DeclareMathOperator{\cH}{H}
\DeclareMathOperator{\dR}{R}
\DeclareMathOperator{\End}{End}
\DeclareMathOperator{\Hasse}{Hasse}
\DeclareMathOperator{\Num}{Num}

\author{James Lee}

% Solarized
% \pagecolor[RGB]{0,20,26}
% \color[RGB]{191,191,191}

% Sephia
% \pagecolor[RGB]{249,239,220}

% Grey
% \pagecolor[RGB]{200, 200, 200}

% Black
% \pagecolor[RGB]{0,0,0}
% \color[RGB]{255,255,255}


\title{Chapter 3, Section 12}

\begin{document}

\maketitle

\begin{enumerate} [label=\textbf{\arabic*.}, leftmargin=0em]

  \item Let $Y$ be a scheme of finite type over an algebraically closed field $k$. Show that the function
        \begin{equation*}
          \varphi(y) = \dim_k(\goth{m}_y/\goth{m}_y^2)
        \end{equation*}
        is upper semicontinuous on the set of closed points of $Y$.

        \begin{proof}
          Apply (12.7.2) to the sheaf of relative differentials $\Omega_{Y/k}$.
          For each closed point $y \in Y$, we have a natural isomorphism (II, 8.7A)
          \begin{equation*}
            \gm_y/\gm_y^2 \cong \Omega_{Y / k} \otimes k(y).
          \end{equation*}
        \end{proof}

  \item Let $\{X_t\}$ be a family of hypersurfaces of the same degree in $\PP_k^n$. Show that for each $i$, the function $h^i(X_t, \fO_{X_t})$ is a constant function of $t$.

        \begin{proof}
          Hypersurfaces of same degree in projective space $\PP^n$ have the same cohomology.
          For any hypersurface $X \subset \PP^n$ of degree $d$, we have an exact sequence
          \[ \begin{tikzcd}
              0 \arrow[r] & \fO_{\PP^n}(-d) \arrow[r] & \fO_{\PP^n} \arrow[r] & \fO_X \arrow[r] & 0,
            \end{tikzcd} \]
          and $\cdh^i(X, \fO_X)$ is independent of the choice of $X$.
        \end{proof}

  \item[\textbf{4.}] Let $Y$ be an integral scheme of finite type over an algebraically closed field $k$. Let $f : X \to Y$ be a flat projective morphism whose fibers are all integral schemes. Let $\fL, \fM$ be invertible sheaves on $X$, and assume for each $y \in Y$ that $\fL_y \cong \fM_y$ on the fiber $X_y$. Then show that there is an invertible sheaf $\fN$ on $Y$ such that $\fL \cong \fM \otimes f^* \fN$.

        \begin{proof}
          If $f$ is flat, then $\fL \otimes \fM^{-1}$ is flat over $Y$.
          For each $y \in Y$, $(\fL \otimes \fM^{-1})_y \cong \fO_{X_y}$ on the fibre $X_y$, and each $X_y$ is integral and projective over $k(y)$, so we have $$\cdh^0(X_y, (\fL \otimes \fM^{-1})_y) = \cdh^0(X_y, \fO_{X_y}) = 1.$$
          Thus, $f_*(\fL \otimes \fM^{-1})$ is locally free of rank one on $Y$ (12.9), say $f^*f_*(\fL \otimes \fM^{-1}) = \fN$ for some $\fN \in \pic{Y}$.
          Applying $f^*$ to both sides, and using the fact that the natural morphism
          \begin{equation*}
            f^*f_*(\fL \otimes \fM^{-1}) \to \fL \otimes \fM^{-1}
          \end{equation*}
          is actually an isomorphism because $f_*(\fL \otimes \fM^{-1})$ is locally free of rank one, we obtain our desired result.
        \end{proof}

  \item[\textbf{5.}] Let $Y$ be an integral scheme of finite type over an algebraically closed field $k$. Let $\mathscr{E}$ be a locally free sheaf on $Y$, and let $X = \PP(\fE)$. Then show that $\pic{X} \cong (\pic{Y}) \times \Z$. This strengthens (II, Ex. 7.9).

        \begin{proof}
          Let $n$ be the rank of $\fE$, $\fO(1)$ the canonical invertible sheaf of $X$, and $\pi : X \to Y$ the natural projection morphism.
          There is an obvious map
          \begin{align*}
            \varphi : (\pic{Y}) \times \Z & \to \pic{X}                      \\
            (\fL, d)                      & \mapsto \fO(d) \otimes \pi^*\fL.
          \end{align*}
          Since $\pi$ is a flat projective morphism, we can apply the results of \S 12.
          In particular, if $\fM$ is an invertible sheaf on $X$, then $\fM_y$ is an invertible sheaf on $X_y \cong \PP_y^n$ for all $y \in Y$, where $\PP_y^n$ is the projective $n$-space over $k(y)$.
          In particular, we can write $\fM_y \cong \fO_{\PP^n_y}(d_y)$ for some $d_y \in \Z$.
          We will show $d_y$ is constant for all $y \in Y$.
          It will follow by (Ex. 12.4) that there exists an invertible sheaf $\fN$ on $Y$ such that $\fM \cong \fO(d) \otimes \pi^*\fN$ for some $d \in \Z$, proving surjectivity of $\varphi$.

          We proceed by showing the minimum value of $d_y$ is obtained when $y = \xi$ is the generic point of $Y$.
          Suppose there exists $y' \in Y$ such that $d_{y'} < d_\xi$.
          After twisting by $\fO(-d_{y'})$, we can assume $d_{y'} = 0$.
          By (12.8), the function
          \begin{equation*}
            \cdh^0(y, \fM) = \dim_{k(y)} \cH^0(\PP^n_y, \fO_{\PP^n_y}(d_y))
          \end{equation*}
          is upper semicontinuous on $Y$ (12.8), and we have the strict inequality
          \begin{equation*}
            \dim_{k(y')} \cH^0(\PP_{y'}^n, \fO_{\PP_{k(y)}}) < \dim_{k(\xi)} \cH^0(\PP_{\xi}^n, \fO_{\PP^n_{\xi}}(d_\xi))
          \end{equation*}
          by the explicit calculations of (5.1).
          But this implies the set $\{ y \in Y | h^0(y, \fM) \geq h^0(\xi, \fM)\}$ is a closed set containing $\xi$ and not containing $y'$, a contradiction.
          Hence, we have $d_\xi \leq d_{y}$ for all $y \in Y$.

          Next, we prove the inequality in the opposite direction.
          By Serre duality (5.1) (7.7), we have
          \begin{align*}
            h^n(y, \fM) & = \dim_{k(y)} \cH^n(\PP_y^n, \fO_{\PP_y^n}(d_y))          \\
                        & = \dim_{k(y)} \cH^n(\PP_y^n, \fO_{\PP_y^n}(-d_y - n - 1)) \\
                        & = h^0(y, \fO(-n-1) \otimes \fM^{-1}),
          \end{align*}
          and $h^n(y, \fM)$ is also an upper semicontinuous on $Y$.
          Suppose there exists $y' \in Y$ such that $d_\xi < d_{y'}$.
          After twisting by $\fO(-d_\xi + n + 1)$, we can assume $d_\xi = n + 1$.
          In particular $d_{y'}$ is a positive integer.
          Again, we have the strict inequality
          \begin{equation*}
            \dim_{k(y')} \cH^0(\PP_{y'}^n, \fO_{\PP_{y'}^n}(-d_{y'} - n - 1)) < \dim_{k(\xi)} \cH^0(\PP_\xi^n, \fO_{\PP_\xi^n}),
          \end{equation*}
          arriving at the same contradiction as above.
          Hence, $d_y$ is constant for all $y \in Y$, and we have shown $\varphi$ is surjective.

          It remains to show $\varphi$ is injective.
          If we have an isomorphism
          \begin{equation*}
            \fO(d) \otimes \pi^*\fL \cong \fO(d') \otimes \pi^* \fL',
          \end{equation*}
          then it is clear by checking on fibres that $d = d'$, so it suffices to check $\pi^* : \pic{Y} \to \pic{X}$ is injective.
          Indeed, we have $\pi_*\fO_X = \fO_Y$ (II, 7.11a), and the projection formula gives
          \begin{equation*}
            \pi_*\pi^*\fN \cong \pi_*\fO_X \otimes \fN \cong \fN
          \end{equation*}
          for all $\fN \in \pic{Y}$, so if $f^*\fN \cong \fO_X$, then $$\fN \cong \pi_*\fO_X \cong \fO_Y.$$
        \end{proof}

  \item[\textbf{6.}] Let $X$ be an integral projective scheme over an algebraically closed field $k$, and assume that $\cH^1(X, \fO_X) = 0$.
        Let $T$ be a connected scheme of finite type over $k$.
        \begin{itemize}
          \item[(a)] If $\fL$ is an invertible sheaf on $X \times T$, show that the invertible sheaves $\fL_t$ on $X = X \times \{ t \}$ are isomorphic, for all closed points $t \in T$.
          \item[(b)] Show that $\pic(X \times T) = \pic{X} \times \pic{T}$.
        \end{itemize}
\end{enumerate}

\begin{proof} $ $ \vspace{0pt}
  \begin{enumerate} [leftmargin=0cm]
    \item[(a)]

    \item[(b)]
  \end{enumerate}
\end{proof}

\end{document}

% \item Let $X_1 \subseteq \PP_k^4$ be the \textit{rational normal quartic curve} (which is the $4$-uple embedding of $\PP^1$ in $\PP^4$). Let $X_0 \subseteq \PP^3$ be a nonsingular rational quartic curve, such as the one in (I, Ex. 3.18b). Use (9.8.3) to construct a flat family $\{X_t\}$ of curves in $\PP^4$, parameterized by $T = \A^1$, with the given fibers $X_1$ and $X_0$ for $t = 1$ and $t = 0$.
% 
% Let $\fI \subseteq \fO_{\PP^4 \times T}$ be the ideal sheaf of the total family $X \subseteq \PP^4 \times T$. Show that $\fI$ is flat over $T$. Then show that
% \begin{equation*}
%   h^0(t, \fI) = \begin{cases}
%     0 & \text{for $t \neq 0$} \\
%     1 & \text{for $t = 0$}
%   \end{cases}
% \end{equation*}
% and also
% \begin{equation*}
%   h^1(t, \fI) = \begin{cases}
%     0 & \text{for $t \neq 0$} \\
%     1 & \text{for $t = 0$}
%   \end{cases}.
% \end{equation*}
% This gives another example of cohomology groups jumping at a special point.
% 
% \begin{proof}
%   Suppose $X_1$ is given by the standard 4-uple embedding $\PP^1$ in $\PP^4$ (I, Ex. 2.12), and $X_0$ is given by (I, Ex. 3.18b).
%   Then, the point $P = (0 : 0 : 1 : 0 : 0)$ does not lie in $X_1$, so by applying (9.8.3)  to the projection $\PP^{4} - P \to \PP^3$ and automorphisms $(x_0 : x_1 : x_2 : x_3 : x_4) \mapsto (x_0 : x_1 : ax_2 : x_3 : x_4)$ for each $a \neq 0$, we have our desired flat family $\{ X_t\}$ of curves in $\PP^4$.
%   To show $\fI$ is flat over $T$, we are just interested in what happens at each affine patch, so we will use affine coordinates $x, y, z, w$ in $\A^4$.
%   Let $X_1$ be given by the parameteric equations
%   \begin{equation*}
%     \begin{cases}
%       x = u^4 \\
%       y = u^3 \\
%       z = u^2 \\
%       w = u
%     \end{cases}.
%   \end{equation*}
%   Now, for any $t \neq 0$, the scheme $X_t$ is given by
%   \begin{equation*}
%     \begin{cases}
%       x = u^4 \\
%       y = u^3 \\
%       z = tu^2 \\
%       w = u
%     \end{cases}.
%   \end{equation*}
%   To get the ideal $I \subseteq k[x, y, z, w]$ of the total family $X$ extended over all of $T$, we eliminate $u$ from the parameteric equations, and make sure $t$ is not a zero divisor in $k[t, x, y, z, w] / I$, so that $X$ will be flat.
%   Note that flatness of $X$ and $\A^4 \times T \cong \A^5$ over $T$ implies flatness of $I$.
%   We find
%   \begin{equation*}
%     I = (t^3y^2 - z^3, t^2x - z^2, tw^2 - z, y - w^3, x - yw).
%   \end{equation*}
%   From this, setting $t = 0$, we obtain the ideal $I_0 \subseteq k[x, y, z, w]$ of $X_0$, which is
%   \begin{equation*}
%     I_0 = (z, y - w^3, x - yw).
%   \end{equation*}
%   So we see that $X_0$ is a schme with support in $\A^3 = \spec{k[x, y, w]}$.
%   It remains to check that $I_0$ gives a nonsingular rational quartic curve in $\A^3$.
%   We proceed by going back to projective space $\PP^3$ and introducing another homogenous coordinate $v$.
%   Consider the surfaces defined by the last two generators of $I_0$
%   \begin{equation*}
%     V : v^2y - w^3 = 0, \quad Q : xv - yw = 0.
%   \end{equation*}
%   The intersection of $V$ and $Q$ should contain another curve of degree 2.
% 
%   By construction, $X_t$ for every $t \neq 0$ is isomorphic to $X_1$, so to compute $h^i(t, \fI)$ for $t \neq 0$, it suffices to compute $\cH^i(\PP^4, \fI_1)$.
%   The exact sequence of $\fO_{\PP^4}$-modules
%   \[ \begin{tikzcd}
%     0 \arrow[r] & \fI_1 \arrow[r] & \fO_{\PP^4} \arrow[r] & \fO_{X_1} \arrow[r] & 0
%   \end{tikzcd} \]
%   induces an exact sequence in cohomology
%   \[ \begin{tikzcd}
% 0 \arrow[r]  & \cH^0(\fI_1) \arrow[r] & \cH^0(\fO_{\PP^4}) \arrow[r] & \cH^0(\fO_{X_1}) \arrow[r] & {} \\
% {} \arrow[r] & \cH^1(\fI_1) \arrow[r] & \cH^1(\fO_{\PP^4})           &                            &   
% \end{tikzcd}\]
%   Since $X_1$ and $\PP^4$ are rational, we have the explicit calculations of (5.1)
%   \begin{equation*}
%     h^0(\fO_{\PP^4}) = 1, \quad h^1(\fO_{X_1}) = 1, \quad h^1(\fO_{\PP^4}) = 0,
%   \end{equation*}
%   and since $\cH^0(\fO_{\PP^4}) \to \cH^1(\fO_{X_1})$ is surjective, by exactness we conclude that $h^i(\fI) = 0$ for $i = 0, 1$.
% \end{proof}
